\chapter*{Abstract}
\addcontentsline{toc}{chapter}{Abstract}

Automated machine learning (AutoML) systems automate parts of the data
science workflow such as model selection, hyperparameter optimisation, and
pipeline construction. Existing AutoML tools are typically either
low-code graphical products that hide their engineering, or Python SDKs
that assume expert operators. Neither approach is well matched to the way
students and researchers actually interact with new datasets --- in an
exploratory, conversational manner --- and neither produces a
reproducible, readable, downloadable trace of the work.

This project introduces \textbf{\projname}, a local, single-user,
\emph{conversational} AutoML platform in which a large language model
(LLM) acts as an autonomous copilot that plans and executes each stage of
the machine-learning workflow by calling well-typed tools. Five in-process
FastMCP microservices --- \svc{mcp-data}, \svc{mcp-eda},
\svc{mcp-modeling}, \svc{mcp-explain}, and \svc{mcp-export} --- expose
eleven tools covering data ingestion, pandera-based schema validation,
exploratory analysis with automatic plot generation, a parallel
scikit-learn/XGBoost/LightGBM ``bake-off'' with cross-validation, an
Optuna hyperparameter sweep, SHAP-based explainability, and
notebook/PDF export. All tools are aggregated into a unified catalog by
an \code{MCPClientPool} that speaks the Model Context Protocol
\citep{mcpSpec2024} directly in-process using \code{FastMCPTransport}.
A Groq-hosted Llama~3.3~70B model \citep{llama3card, groq} drives the
tool-calling loop, and a thin Streamlit front-end \citep{streamlit}
provides upload, chat, and artefact-browsing interfaces.

The design emphasises three principles: (i)~the LLM is only used for
planning and language --- all numerical work is done by deterministic,
inspectable Python code; (ii)~every artefact (chart, model, notebook,
report) is materialised to a local, indexed filesystem and can be
opened, downloaded, or re-executed without the platform running;
and (iii)~each pipeline stage is an independently swappable MCP server,
so new tools can be added without modifying the orchestrator or the UI.

We describe the system architecture, MCP integration, tool schemas,
orchestration loop, artefact registry, and Streamlit interface. We
present the intended experimental protocol, discuss the limitations of
the current implementation (single user, local storage, hallucination
risk in tool selection, dependency footprint), and outline realistic
future extensions such as multi-agent workflows, distributed training,
experiment tracking, and a model registry.

\vspace{0.5cm}
\noindent\textbf{Keywords:} Automated Machine Learning, Model Context
Protocol, Large Language Model Agents, Tool Use, Streamlit, Groq,
scikit-learn, XGBoost, LightGBM, Optuna, SHAP, Reproducible Machine
Learning.

\clearpage
